% =============================================================================
% 01-fenomenologia.tex - PARTE 1, o que estamos tentando entender.
% Frames 04 a 08.
% =============================================================================

% ----------------------------------------------------------------- frame 04 --
% Primeiro frame de conteudo. Declara o recorte em uma linha e traz o mapa
% fenomenologico NAO linear, com relacoes bidirecionais.
\begin{frame}{\ftit{\faSitemap}{O que estamos tentando entender}}
\CRar{2.5}
{\footnotesize\sffamily\color{campografite}%
A prele\c{c}\~ao paralela do m\'odulo cobre o diagn\'ostico do epis\'odio
depressivo; esta ATP cobre a conduta e trata a suicidalidade como fen\^omeno
transdiagn\'ostico.\par}

\CRvg
\begin{center}
\begin{tikzpicture}[node distance=0.42cm and 0.60cm]
  \node[crnode={campoazul}{2.30cm}, font=\scriptsize]                 (pass) {idea\c{c}\~ao\\passiva};
  \node[crnode={campoazul}{2.30cm}, font=\scriptsize, right=of pass]  (ativ) {idea\c{c}\~ao\\ativa};
  \node[crnode={campoazul}{2.30cm}, font=\scriptsize, right=of ativ]  (inte) {inten\c{c}\~ao};
  \node[crnode={campoambar}{2.30cm}, font=\scriptsize, below=0.72cm of pass] (plan) {planejamento};
  \node[crnode={campoambar}{2.30cm}, font=\scriptsize, right=of plan] (prep) {comportamento\\preparat\'orio};
  \node[crnode={campobordo}{2.30cm}, font=\scriptsize, right=of prep] (tent) {tentativa};

  \draw[musialBi] (pass.east) -- (ativ.west);
  \draw[musialBi] (ativ.east) -- (inte.west);
  \draw[musialBi] (plan.east) -- (prep.west);
  \draw[musialBi] (prep.east) -- (tent.west);
  \draw[musialBi] (inte.south) -- (prep.north);
  \draw[musialBi] (ativ.south) -- (plan.north);
  \draw[musialBi, campobordo!80!black] (tent.north) to[out=60,in=0,looseness=1.5] (inte.east);
  \draw[musialBi] (plan.west) to[out=180,in=180,looseness=1.3] (pass.west);

  \node[crnode={campografite}{3.05cm}, font=\scriptsize, below=0.62cm of prep]
    (auto) {autoles\~ao sem inten\c{c}\~ao suicida\\[1pt]\scriptsize eixo paralelo};
  \draw[musialThin, dash pattern=on 2.4pt off 2pt] (auto.north) -- (prep.south);
\end{tikzpicture}
\end{center}

\mensagem{Esses fen\^omenos se relacionam nos dois sentidos. N\~ao s\~ao uma
escada obrigat\'oria.}
\end{frame}

% ----------------------------------------------------------------- frame 05 --
% A pirâmide de progressao como objeto de critica, e o eixo paralelo da autolesao.
\begin{frame}{\ftit{\faBan}{Duas leituras que atrapalham}}
\CRar{6}
\begin{columns}[onlytextwidth, t]
\begin{column}{0.47\textwidth}
\begin{caixaagir}{A escada que n\~ao existe}
Ningu\'em precisa passar por todos os degraus, na ordem, para chegar \`a
tentativa. Uma parte das tentativas ocorre em quem nunca planejou, e uma parte
de quem planeja durante anos nunca tenta.\par
\agir{Consequ\^encia:} a aus\^encia de um degrau anterior n\~ao tranquiliza.
\end{caixaagir}
\end{column}
\calha
\begin{column}{0.47\textwidth}
\begin{caixaavaliar}{A autoles\~ao no eixo ao lado}
A autoles\~ao sem inten\c{c}\~ao suicida tem fun\c{c}\~ao de regula\c{c}\~ao do
sofrimento, e n\~ao de morte. N\~ao \'e tentativa.\par
\avaliar{E ainda assim} \'e um dos marcadores mais consistentes de risco
futuro. As duas coisas s\~ao verdadeiras ao mesmo tempo.
\end{caixaavaliar}
\end{column}
\end{columns}

\CRvg
\begin{center}
\begin{tikzpicture}
  \node[crnode={campografite}{9.6cm}, font=\scriptsize] {%
    Dizer \fala{ele s\'o tem idea\c{c}\~ao passiva} e encerrar a\'i \'e usar um
    degrau como se fosse um progn\'ostico.};
\end{tikzpicture}
\end{center}
\end{frame}

% ----------------------------------------------------------------- frame 06 --
% Magnitude e tendencia. pgfplots. Apenas os dois pontos documentados.
\begin{frame}{\ftit{\faChartLine}{Magnitude e tend\^encia no Brasil}}
\CRar{4.6}
\begin{columns}[onlytextwidth, T]
\begin{column}{0.50\textwidth}
\begin{center}
\begin{tikzpicture}
\begin{axis}[
  width=5.4cm, height=4.2cm,
  ymin=0, ymax=9, xmin=2008.6, xmax=2022.4,
  xtick={2010,2021}, xticklabels={2010,2021},
  ytick={0,2,4,6,8},
  ylabel={\scriptsize \'obitos por 100 mil},
  y label style={font=\scriptsize, color=campografite},
  tick label style={font=\scriptsize, color=campografite},
  axis line style={campografite!55},
  grid=major, grid style={campografite!14, line width=0.3pt},
  every axis plot/.append style={line width=1.1pt}]
  \addplot[campoazul!85!black, mark=*, mark size=2.3pt,
           mark options={fill=campoazul!25, draw=campoazul!85!black}]
    coordinates {(2010,5.2) (2021,7.5)};
  \node[font=\scriptsize\sffamily, color=campoazul!85!black, anchor=south east]
    at (axis cs:2010,5.2) {5,2};
  \node[font=\scriptsize\sffamily, color=campoazul!85!black, anchor=south east]
    at (axis cs:2021,7.5) {7,5};
\end{axis}
\end{tikzpicture}
\end{center}
{\scriptsize\sffamily\color{campografite}\hspace{0.3cm}Dois pontos documentados.
A s\'erie intermedi\'aria n\~ao foi reconstru\'ida.\par}
\end{column}
\calha
\begin{column}{0.44\textwidth}
\begin{caixaavaliar}{2021, no Brasil}
\begin{itemc}
\item \avaliar{15.507} mortes por suic\'idio
\item \avaliar{114.159} notifica\c{c}\~oes de les\~ao autoprovocada
\item taxa de \avaliar{7,5} por 100 mil, \mudanca{42\% acima} da taxa de 2010
\end{itemc}
\end{caixaavaliar}
\CRvg
\begin{caixamudanca}{O que essa raz\~ao diz}
Para cada morte registrada h\'a cerca de sete notifica\c{c}\~oes de les\~ao
autoprovocada. Essas pessoas passam pelo servi\c{c}o.
\end{caixamudanca}
\end{column}
\end{columns}
\fonte{Minist\'erio da Sa\'ude, Boletim Epidemiol\'ogico v.\,55, n.\,4, 2024,
com dados do SIM e do SINAN.}
\end{frame}

% ----------------------------------------------------------------- frame 07 --
% Faixas etarias e sexo. Barras agrupadas, dado numericamente publicado.
\begin{frame}{\ftit{\faUsers}{Quem morre, por faixa et\'aria e sexo}}
\CRar{6}
\begin{center}
\begin{tikzpicture}
\begin{axis}[
  width=10.6cm, height=5.0cm,
  ybar=1.4pt, bar width=8.5pt,
  ymin=0, ymax=23, ytick={0,5,10,15,20},
  symbolic x coords={5 a 14,15 a 19,20 a 29,30 a 49,50 a 69,70 ou mais},
  xtick=data,
  ylabel={\scriptsize \'obitos por 100 mil, 2021},
  y label style={font=\scriptsize, color=campografite},
  tick label style={font=\scriptsize, color=campografite},
  axis line style={campografite!55},
  ymajorgrids, grid style={campografite!14, line width=0.3pt},
  legend style={font=\scriptsize\sffamily, draw=campografite!35, at={(0.015,0.98)},
                anchor=north west, legend columns=2, column sep=5pt},
  nodes near coords, nodes near coords style={font=\tiny\sffamily,
                                             color=campografite},
  every node near coord/.append style={/pgf/number format/.cd, fixed,
                                       precision=1, use comma}]
  \addplot[draw=campoazul!85!black, fill=campoazul!30,
           every node near coord/.append style={/tikz/.cd, yshift=5.5pt, xshift=-2.6pt}]
    coordinates {(5 a 14,0.7) (15 a 19,9.3) (20 a 29,14.6) (30 a 49,14.9)
                 (50 a 69,15.4) (70 ou mais,18.1)};
  \addplot[draw=camporoxo!80!black, fill=camporoxo!22,
           every node near coord/.append style={/tikz/.cd, xshift=2.6pt}]
    coordinates {(5 a 14,0.9) (15 a 19,4.5) (20 a 29,3.9) (30 a 49,3.8)
                 (50 a 69,3.8) (70 ou mais,2.9)};
  \legend{homens, mulheres}
\end{axis}
\end{tikzpicture}
\end{center}

\mensagem{A taxa mais alta est\'a nos homens com 70 anos ou mais. \'E a faixa que
o servi\c{c}o menos pergunta.}
\fonte{Minist\'erio da Sa\'ude, Boletim Epidemiol\'ogico v.\,55, n.\,4, 2024.}
\end{frame}

% ----------------------------------------------------------------- frame 08 --
% Contato previo com servicos. O argumento que sustenta toda a aula.
\begin{frame}{\ftit{\faUserMd}{Eles passam por n\'os antes}}
\CRar{6}
\begin{center}
\begin{tikzpicture}
\begin{axis}[
  width=9.55cm, height=4.10cm,
  xbar, bar width=11pt,
  xmin=0, xmax=100,
  symbolic y coords={smm, apm, sma, apa},
  ytick=data,
  yticklabels={{sa\'ude mental, no m\^es},
               {aten\c{c}\~ao prim\'aria, no m\^es},
               {sa\'ude mental, no ano},
               {aten\c{c}\~ao prim\'aria, no ano}},
  yticklabel style={font=\scriptsize, color=campografite, align=right},
  xlabel={\scriptsize \% das pessoas que morreram por suic\'idio},
  x label style={font=\scriptsize, color=campografite},
  tick label style={font=\scriptsize, color=campografite},
  axis line style={campografite!55},
  enlarge y limits=0.20,
  xmajorgrids, grid style={campografite!14, line width=0.3pt},
  nodes near coords, nodes near coords style={font=\tiny\sffamily,
                                              color=campografite, anchor=west,
                                              xshift=1.5pt},
  point meta=explicit symbolic]
  \addplot[draw=campoazul!85!black, fill=campoazul!30]
    coordinates {(20,smm)[20\%] (45,apm)[45\%]
                 (33,sma)[33\%] (75,apa)[75\%]};
\end{axis}
\end{tikzpicture}
\end{center}

\begin{caixaavaliar}{O que isso obriga}
A pessoa que morre por suic\'idio costuma ter estado diante de um m\'edico pouco
antes, e quase sempre de um m\'edico que n\~ao era psiquiatra. Em idosos, o
contato com a aten\c{c}\~ao prim\'aria no m\^es anterior \'e ainda mais
frequente do que em adultos jovens.
\end{caixaavaliar}

\fonte{Luoma, Martin e Pearson, \emph{Am J Psychiatry} 2002, revis\~ao de 40
estudos. \seloevidencia}
\end{frame}
